\documentclass[11pt]{article}
\usepackage[margin=0.9in]{geometry}
\usepackage{graphicx,amsmath,amssymb,booktabs,hyperref,xcolor,tabularx,float}
\graphicspath{{./}}
\hypersetup{colorlinks=true,linkcolor=blue,urlcolor=blue}
\setlength{\parskip}{6pt}

\title{$(L_{\rm priv}, L_{\rm pub})$ coordinate transformation:\\
       test of the brainstorm hypothesis\\
       \large{Addendum to the K=3 study summary}}
\author{Fixed-Point Factory}
\date{June 2026}

\begin{document}
\maketitle

\begin{abstract}
\noindent
The hypothesis: transform $u_k$ to ``private-only log-odds''
$L_{\rm priv}(u_k) = \tau u_k$ and $p$ to ``public-only log-odds''
$L_{\rm pub}(p) = \log[f_1(p)/f_0(p)]$. Bayes then suggests
$\operatorname{logit}\mu(p, u_k) = L_{\rm priv} + L_{\rm pub} + \delta$
with $\delta$ a small smooth $K=3$ conditional-overlap correction.
If so, fitting Cheby in $(L_{\rm priv}, L_{\rm pub})$ coordinates
would converge spectrally with very few coefficients.

\smallskip
\textbf{Verdict: partially confirmed.} The additive decomposition
correctly identifies the dominant structure --- $\delta_{\rm RMS}$ is
$\sim$50\% smaller than $L_{\rm RMS}$ at every $(\gamma, \tau)$ tested.
\textbf{But the Cheby coefficient decay rate of $\delta$ is essentially
identical to that of $L$:} the cusps in $\mu(p, u_k)$ (induced by
critical points of $P$, see main report \S 18) are preserved unchanged
in $\delta$, so the smoothness barrier survives the coordinate
transformation. The SVD rank of $\delta$ stays close to full ($\sim G$)
just like $L$.

\smallskip
\textbf{Useful upshot.} The decomposition is useful for storage and
interpretation (we now have explicit closed-form expressions for the
dominant additive structure) but does not unlock spectral convergence.
The remaining $\sim$50\% of variation is intrinsically high-rank
cuspy structure that requires \emph{piecewise/multi-element} treatment
to resolve --- consistent with the main report's earlier conclusion
about the smoothness barrier.

\smallskip
\textbf{Concrete economic finding.} The fitted $b(u_k)$ effect is
$\sim 10{\times}$ smaller than the textbook closed-form $\tau \cdot u_k$
prediction at $\tau = 1$, $\gamma = 100$. This is itself meaningful:
at high $\tau$, the price already aggregates the agent's signal so
well that $u_k$'s marginal information contribution to $\mu(p, u_k)$
is much smaller than the linear-Gaussian limit suggests.
\end{abstract}

\section{Test methodology}

We loaded the strict-$h{=}0$ fixed point at four $(\gamma, \tau)$ cells
saved in \texttt{dd\_k3\_strict\_fp\_g\{10,100,1000\}\_t\{0.2,1.0\}.npz}
and extracted the lookup matrix
$\mu_{\rm strict}(p_i, u_{k,j})$ on the $(G_p, G) = (121, 11)$ grid.
We then computed
\[
  L(p_i, u_{k,j}) \;:=\; \operatorname{logit}\mu_{\rm strict}(p_i, u_{k,j}),
\]
and tested two decompositions:

\begin{itemize}\itemsep1pt
  \item \textbf{ANOVA additive:} $L = \bar L + a(p) + b(u_k) + \delta$
        with $\bar a = \bar b = 0$, minimizing $\|\delta\|_2$ in
        closed form via row/column means.
  \item \textbf{Cheby-h additive:} same form but $b(u_k)$ constrained
        to a degree-8 Cheby polynomial in $u_k$; $a(p)$ free per node.
        Solved by alternating LSQ.
\end{itemize}

For each decomposition we computed the Cheby coefficient tail-norm
profile $\|c\|_2(\text{order} \ge N)$, the SVD spectrum of $\delta$,
and the rank-$k$ truncation error.

\section{Result 1: bulk absorbed, cusps preserved}

\begin{figure}[H]\centering
  \includegraphics[width=\linewidth]{lpub_figs/fig2_heat.png}
  \caption{Heatmaps of $L(p, u_k)$, $\delta_{\rm ANOVA}$, and
    $\delta_{\rm Cheby-h}$ across four cells. The smooth bulk
    (visible in $L$ as a vertical gradient) is removed in $\delta$,
    but the sharp horizontal bands at the critical-value $p$'s
    (cusps) are preserved unchanged. The two decompositions give
    visually identical residuals.}
\end{figure}

The key visual: in the $\delta$ panels (middle and right columns),
the smooth blue-to-red gradient that dominates $L$ is gone, but the
sharp horizontal bands at $p \approx 0.34$ and $p \approx 0.66$
(at $\tau = 1$, $\gamma = 100$) survive. Those bands are where
$|\nabla P|$ vanishes at critical points of $P$; they generate the
$\log$-cusps in the co-area integral.

\section{Result 2: identical Cheby decay rates}

\begin{figure}[H]\centering
  \includegraphics[width=\linewidth]{lpub_figs/fig1_tail_decay.png}
  \caption{Cumulative Cheby-coefficient tail norm:
    $\bigl(\sum_{i+j \ge N} c_{ij}^2\bigr)^{1/2}$ vs $N$, for the
    direct Cheby fit of $L$ (blue), of $\delta_{\rm ANOVA}$ (red),
    and of $\delta_{\rm Cheby-h}$ (green). Across all four cells
    the three curves are essentially indistinguishable. The
    coordinate transformation does not improve the per-coefficient
    convergence rate.}
\end{figure}

If the cusps had been absorbed into the coordinate change, the red
curve would lie a decade or two below the blue. They don't: the
slow algebraic decay is preserved.

\section{Result 3: SVD rank of $\delta$ matches $L$}

\begin{figure}[H]\centering
  \includegraphics[width=\linewidth]{lpub_figs/fig4_svd.png}
  \caption{Singular values of $\delta$ vs $L$ (both normalized to
    $\sigma_1$). The two spectra overlap through $k \approx 9$. The
    cusps make $L$ effectively full-rank ($G = 11$), and $\delta$
    inherits the same rank. Rank-1 truncation captures only 30--45\%
    of the Frobenius norm of $\delta$; rank-5 captures 65--80\%; you
    need essentially all of the singular vectors to reach
    machine $\varepsilon$.}
\end{figure}

\begin{center}\begin{tabular}{lrrrr}\toprule
$(\gamma, \tau)$ & rank-1 & rank-2 & rank-3 & rank-5 \\
\midrule
$(100, 0.2)$  & 0.446 & 0.556 & 0.671 & 0.878 \\
$(100, 1.0)$  & 0.348 & 0.515 & 0.617 & 0.770 \\
$(10, 1.0)$   & 0.331 & 0.489 & 0.576 & 0.757 \\
$(1000, 1.0)$ & 0.306 & 0.452 & 0.527 & 0.665 \\
\bottomrule\end{tabular}\end{center}

Fraction of Frobenius norm captured by rank-$k$ truncation of
$\delta_{\rm ANOVA}$. A clean low-rank correction would need rank 1
or 2 to capture $>$95\%; we get $<$45\%. So a multiplicative-rank-1
extension like $L \approx a(u_k) + b(p) + c(u_k) \cdot d(p)$ would
not close the gap either.

\section{Result 4: the empirical $b(u_k)$ effect is much smaller than $\tau u_k$}

\begin{figure}[H]\centering
  \includegraphics[width=0.85\linewidth]{lpub_figs/fig3_effects.png}
  \caption{Row-effect $a(p)$ (left column) and column-effect $b(u_k)$
    (right column). The dotted black line in the right panels is the
    naive textbook prediction $L_{\rm priv}(u_k) = \tau u_k$. The
    fitted effect (solid green) is $\sim 10{\times}$ smaller in
    magnitude. The Cheby-h fit recovers essentially the same shape as
    the ANOVA effect.}
\end{figure}

At $(\gamma, \tau) = (100, 1)$: the textbook $L_{\rm priv}(u_k)$
ranges $[-2.33, +2.33]$ (since $u_k \in [-2.33, +2.33]$ and $\tau = 1$),
but the empirical $b(u_k)$ ranges only $[-0.22, +0.22]$. The ``naive
private log-odds'' over-estimates the $u_k$ contribution to $\mu$ by
an order of magnitude.

\textbf{Economic interpretation.} At high $\tau$ the equilibrium
price aggregates the signals so effectively that, conditional on $p$,
each agent's own $u_k$ adds little marginal information. The
posterior $\mu(p, u_k)$ is dominated by $p$. This is consistent with
the convergence of K-agent CRRA REE to the perfect-information limit
as $K \to \infty$: $K = 3$ already does most of the aggregation work.

\section{Why the brainstorm didn't quite work}

The hypothesis was: in the right coordinates, the lookup is affine +
a small smooth correction, so Cheby converges spectrally on the
correction.

What's actually true: in the right coordinates, the lookup is affine
$+$ a smooth correction $+$ \emph{the cusps}. The cusps are
preserved by any global linear change of coordinates --- they live on
the algebraic variety where $|\nabla P| = 0$, which is intrinsic to
$P$ and survives any pre-composition. To remove them you have to
either (i) detect them and use piecewise representations (multi-element
spectral), or (ii) use a representation like PCHIP that's natively
$C^1$-not-$C^\infty$ (what Option B+ does).

The decomposition is still useful for two things:
\begin{enumerate}\itemsep1pt
  \item \textbf{Storage / serialization.} Store the additive parts
        $a(p)$ and $b(u_k)$ as 1D arrays (much smaller than the 2D
        $\mu(p, u_k)$ table) and reconstruct on the fly; the
        residual $\delta$ is still 2D but with lower magnitude, so
        a coarser PCHIP grid for $\delta$ might suffice.
  \item \textbf{Diagnostic / interpretation.} The empirical
        $b(u_k) \ll \tau u_k$ is an interesting finding about K=3
        information aggregation; useful for the economics writeup
        even if not for the numerics.
\end{enumerate}

\section{Implication for Method C v3 design}

The non-additive correction we were planning to add as
$P = \sigma(h(u_1) + h(u_2) + h(u_3) + \kappa \cdot g(u_1, u_2, u_3))$
is for $P$, not $\mu$. The cusps in $\mu$ would still emerge under
v3 because they come from $|\nabla P|^{-1}$ in the co-area integral
--- they're a property of the level-set geometry of $P$, regardless
of how smooth $P$ itself is. So Method C v3 would lower
$\|F\|_\infty$ at $\tau = 1$ (capturing more of the additive
deviation from the linear-Gaussian limit) but would not eliminate
the cusps in the resulting $\mu(p, u_k)$.

This suggests two independent directions for the next step:
\begin{enumerate}\itemsep1pt
  \item \textbf{Method C v3 for $P$:} smooth, finite-parameter
        ansatz, drives $F$ residual toward zero across the
        $(\gamma, \tau)$ plane. Useful for fast continuation
        studies; doesn't help lookup accuracy beyond what
        Option B+ already does.
  \item \textbf{Cusp-aware $\mu$ representation:} keep using
        Option B+ (PCHIP density) which natively handles cusps, or
        extend to piecewise Cheby with critical-value knots if
        machine-$\varepsilon$ lookups are needed.
\end{enumerate}

These are complementary: (1) gives a fast smooth $P$ for continuation
and the bulk of equilibrium analysis; (2) gives accurate $\mu$ where
needed.

\section{Reproducibility}

Code: \texttt{/tmp/dd\_k3\_lpub\_lpriv\_v2.py} (and SVD diagnostic
\texttt{/tmp/dd\_k3\_delta\_svd.py}). To re-run, you need the saved
strict-$h{=}0$ FPs in
\texttt{projects/REZN/solved\_fixed\_points/dd\_k3\_overnight/dd\_k3\_strict\_fp\_*.npz};
all four-cell results take $\sim 5$\,seconds.

Output: \texttt{lpub\_lpriv\_v2/figs/} and \texttt{results.json}.

\end{document}
